% Partie : incomplétude
% Chapitre : représentabilité dans Q
% Section : C

\documentclass[../../../include/open-logic-section]{subfiles}

\begin{document}

\olfileid[fr]{inc}{req}{cee}
\olsection{Les fonctions de $C$}

Soit $C$ le plus petit ensemble de fonctions contenant :
\begin{enumerate}
\item $0$ ;
\item le successeur ;
\item l'addition ;
\item la multiplication ;
\item les projections ;
\item la fonction caractéristique de l'égalité, $\Char{=}$ ;
\end{enumerate}
et clos par :
\begin{enumerate}
\item composition ;
\item recherche non bornée appliquée aux fonctions régulières.
\end{enumerate}
Cette dernière restriction signifie que l'opération $\mu$
ne peut être employée que lorsque la fonction obtenue est
totale. Comparons cette définition à celle des fonctions
\emph{récursives générales} : nous avons ajouté l'addition,
la multiplication et $\Char{=}$, mais supprimé la récurrence
primitive.

Toute fonction de $C$ est manifestement récursive, puisque
l'addition, la multiplication et $\Char{=}$ le sont.
Nous montrerons que la réciproque est également vraie :
les autres fonctions et opérations de $C$ permettent
d'effectuer la récurrence primitive.

\end{document}
